% Options for packages loaded elsewhere
\PassOptionsToPackage{unicode}{hyperref}
\PassOptionsToPackage{hyphens}{url}
%
\documentclass[
]{article}
\usepackage{amsmath,amssymb}
\usepackage{iftex}
\ifPDFTeX
  \usepackage[T1]{fontenc}
  \usepackage[utf8]{inputenc}
  \usepackage{textcomp} % provide euro and other symbols
\else % if luatex or xetex
  \usepackage{unicode-math} % this also loads fontspec
  \defaultfontfeatures{Scale=MatchLowercase}
  \defaultfontfeatures[\rmfamily]{Ligatures=TeX,Scale=1}
\fi
\usepackage{lmodern}
\ifPDFTeX\else
  % xetex/luatex font selection
\fi
% Use upquote if available, for straight quotes in verbatim environments
\IfFileExists{upquote.sty}{\usepackage{upquote}}{}
\IfFileExists{microtype.sty}{% use microtype if available
  \usepackage[]{microtype}
  \UseMicrotypeSet[protrusion]{basicmath} % disable protrusion for tt fonts
}{}
\makeatletter
\@ifundefined{KOMAClassName}{% if non-KOMA class
  \IfFileExists{parskip.sty}{%
    \usepackage{parskip}
  }{% else
    \setlength{\parindent}{0pt}
    \setlength{\parskip}{6pt plus 2pt minus 1pt}}
}{% if KOMA class
  \KOMAoptions{parskip=half}}
\makeatother
\usepackage{xcolor}
\usepackage[margin=1in]{geometry}
\usepackage{color}
\usepackage{fancyvrb}
\newcommand{\VerbBar}{|}
\newcommand{\VERB}{\Verb[commandchars=\\\{\}]}
\DefineVerbatimEnvironment{Highlighting}{Verbatim}{commandchars=\\\{\}}
% Add ',fontsize=\small' for more characters per line
\usepackage{framed}
\definecolor{shadecolor}{RGB}{248,248,248}
\newenvironment{Shaded}{\begin{snugshade}}{\end{snugshade}}
\newcommand{\AlertTok}[1]{\textcolor[rgb]{0.94,0.16,0.16}{#1}}
\newcommand{\AnnotationTok}[1]{\textcolor[rgb]{0.56,0.35,0.01}{\textbf{\textit{#1}}}}
\newcommand{\AttributeTok}[1]{\textcolor[rgb]{0.13,0.29,0.53}{#1}}
\newcommand{\BaseNTok}[1]{\textcolor[rgb]{0.00,0.00,0.81}{#1}}
\newcommand{\BuiltInTok}[1]{#1}
\newcommand{\CharTok}[1]{\textcolor[rgb]{0.31,0.60,0.02}{#1}}
\newcommand{\CommentTok}[1]{\textcolor[rgb]{0.56,0.35,0.01}{\textit{#1}}}
\newcommand{\CommentVarTok}[1]{\textcolor[rgb]{0.56,0.35,0.01}{\textbf{\textit{#1}}}}
\newcommand{\ConstantTok}[1]{\textcolor[rgb]{0.56,0.35,0.01}{#1}}
\newcommand{\ControlFlowTok}[1]{\textcolor[rgb]{0.13,0.29,0.53}{\textbf{#1}}}
\newcommand{\DataTypeTok}[1]{\textcolor[rgb]{0.13,0.29,0.53}{#1}}
\newcommand{\DecValTok}[1]{\textcolor[rgb]{0.00,0.00,0.81}{#1}}
\newcommand{\DocumentationTok}[1]{\textcolor[rgb]{0.56,0.35,0.01}{\textbf{\textit{#1}}}}
\newcommand{\ErrorTok}[1]{\textcolor[rgb]{0.64,0.00,0.00}{\textbf{#1}}}
\newcommand{\ExtensionTok}[1]{#1}
\newcommand{\FloatTok}[1]{\textcolor[rgb]{0.00,0.00,0.81}{#1}}
\newcommand{\FunctionTok}[1]{\textcolor[rgb]{0.13,0.29,0.53}{\textbf{#1}}}
\newcommand{\ImportTok}[1]{#1}
\newcommand{\InformationTok}[1]{\textcolor[rgb]{0.56,0.35,0.01}{\textbf{\textit{#1}}}}
\newcommand{\KeywordTok}[1]{\textcolor[rgb]{0.13,0.29,0.53}{\textbf{#1}}}
\newcommand{\NormalTok}[1]{#1}
\newcommand{\OperatorTok}[1]{\textcolor[rgb]{0.81,0.36,0.00}{\textbf{#1}}}
\newcommand{\OtherTok}[1]{\textcolor[rgb]{0.56,0.35,0.01}{#1}}
\newcommand{\PreprocessorTok}[1]{\textcolor[rgb]{0.56,0.35,0.01}{\textit{#1}}}
\newcommand{\RegionMarkerTok}[1]{#1}
\newcommand{\SpecialCharTok}[1]{\textcolor[rgb]{0.81,0.36,0.00}{\textbf{#1}}}
\newcommand{\SpecialStringTok}[1]{\textcolor[rgb]{0.31,0.60,0.02}{#1}}
\newcommand{\StringTok}[1]{\textcolor[rgb]{0.31,0.60,0.02}{#1}}
\newcommand{\VariableTok}[1]{\textcolor[rgb]{0.00,0.00,0.00}{#1}}
\newcommand{\VerbatimStringTok}[1]{\textcolor[rgb]{0.31,0.60,0.02}{#1}}
\newcommand{\WarningTok}[1]{\textcolor[rgb]{0.56,0.35,0.01}{\textbf{\textit{#1}}}}
\usepackage{graphicx}
\makeatletter
\newsavebox\pandoc@box
\newcommand*\pandocbounded[1]{% scales image to fit in text height/width
  \sbox\pandoc@box{#1}%
  \Gscale@div\@tempa{\textheight}{\dimexpr\ht\pandoc@box+\dp\pandoc@box\relax}%
  \Gscale@div\@tempb{\linewidth}{\wd\pandoc@box}%
  \ifdim\@tempb\p@<\@tempa\p@\let\@tempa\@tempb\fi% select the smaller of both
  \ifdim\@tempa\p@<\p@\scalebox{\@tempa}{\usebox\pandoc@box}%
  \else\usebox{\pandoc@box}%
  \fi%
}
% Set default figure placement to htbp
\def\fps@figure{htbp}
\makeatother
\setlength{\emergencystretch}{3em} % prevent overfull lines
\providecommand{\tightlist}{%
  \setlength{\itemsep}{0pt}\setlength{\parskip}{0pt}}
\setcounter{secnumdepth}{-\maxdimen} % remove section numbering
\usepackage{bookmark}
\IfFileExists{xurl.sty}{\usepackage{xurl}}{} % add URL line breaks if available
\urlstyle{same}
\hypersetup{
  pdftitle={Dever de classe 01 - Dados Hubble},
  pdfauthor={Cecilia Lodron Gonzaga e Francielle Rezende},
  hidelinks,
  pdfcreator={LaTeX via pandoc}}

\title{Dever de classe 01 - Dados Hubble}
\author{Cecilia Lodron Gonzaga e Francielle Rezende}
\date{2026-09-04}

\begin{document}
\maketitle

\section{Resumo}\label{resumo}

Neste relatório, iremos verificar os dados do satélite Hubble e
verificarmos a linearidade dos dados entre a distância e a velocidade.
Primeiramente, verificaremos o comportamento dos dados logo após
verificar se podemos encaixar uma regressão linear entre as variáveis.

\section{Introdução}\label{introduuxe7uxe3o}

Em 1929, os astrônomo Edwin Hubble, trabalhando no observatório Mount
Wilson, realizou um estudo em que foram analisadas as velocidades e
distâncias entre as galáxias. Hubble utilizou as chamadas estrelas
Cafeidas para medir as distâncias das várias galáxias. Esses dados foram
publicados e, segundo Machado, foram considerados um grande marco na
cosmologia.\\
Nesse trabalho iremos utilizar os dados obtidos por Hubble para avaliar
a relação linear entre a velocidade e a distância das galáxias.

\begin{Shaded}
\begin{Highlighting}[]
\FunctionTok{Sys.setenv}\NormalTok{(}\AttributeTok{LANG =} \StringTok{"en"}\NormalTok{)}
\FunctionTok{rm}\NormalTok{(}\AttributeTok{list =} \FunctionTok{ls}\NormalTok{(}\AttributeTok{all =} \ConstantTok{TRUE}\NormalTok{))}
\FunctionTok{setwd}\NormalTok{(}\StringTok{"C:/Users/cecil/OneDrive/Documentos"}\NormalTok{)}
\NormalTok{dados }\OtherTok{\textless{}{-}} \FunctionTok{read\_excel}\NormalTok{(}\AttributeTok{path =} \StringTok{"1929\_Hubble.xlsx"}\NormalTok{, }\AttributeTok{sheet =} \StringTok{"Data"}\NormalTok{)}

\NormalTok{dados[}\DecValTok{1}\SpecialCharTok{:}\DecValTok{15}\NormalTok{,}\FunctionTok{c}\NormalTok{(}\DecValTok{3}\NormalTok{,}\DecValTok{5}\NormalTok{)]}
\end{Highlighting}
\end{Shaded}

\begin{verbatim}
## # A tibble: 15 x 2
##    v_km_s r_Mpc
##     <dbl> <chr>
##  1    170 0.032
##  2    290 0.034
##  3   -130 0.214
##  4    -70 0.263
##  5   -185 0.275
##  6   -220 0.275
##  7    200 0.45 
##  8    290 0.5  
##  9    270 0.5  
## 10    200 0.63 
## 11    300 0.8  
## 12    -30 0.9  
## 13    650 0.9  
## 14    150 0.9  
## 15    500 0.9
\end{verbatim}

\begin{Shaded}
\begin{Highlighting}[]
\FunctionTok{par}\NormalTok{(}\AttributeTok{mfrow =} \FunctionTok{c}\NormalTok{(}\DecValTok{1}\NormalTok{,}\DecValTok{2}\NormalTok{))}

\CommentTok{\#guardando as colunas em variáveis}
\NormalTok{velocidade }\OtherTok{\textless{}{-}}\NormalTok{ dados}\SpecialCharTok{$}\NormalTok{v\_km\_s}
\NormalTok{distancia }\OtherTok{\textless{}{-}} \FunctionTok{as.numeric}\NormalTok{(dados}\SpecialCharTok{$}\NormalTok{r\_Mpc)}

\CommentTok{\#Resumo dos dados}
\FunctionTok{summary}\NormalTok{(velocidade)}
\end{Highlighting}
\end{Shaded}

\begin{verbatim}
##    Min. 1st Qu.  Median    Mean 3rd Qu.    Max. 
##  -220.0   290.0   590.0   552.6   800.0  1800.0
\end{verbatim}

\begin{Shaded}
\begin{Highlighting}[]
\FunctionTok{summary}\NormalTok{(distancia)}
\end{Highlighting}
\end{Shaded}

\begin{verbatim}
##    Min. 1st Qu.  Median    Mean 3rd Qu.    Max.    NA's 
##   0.032   0.670   1.100   1.213   1.730   3.450       1
\end{verbatim}

\begin{Shaded}
\begin{Highlighting}[]
\CommentTok{\#Boxplot da velocidade}
\FunctionTok{boxplot}\NormalTok{(cars}\SpecialCharTok{$}\NormalTok{speed,}
        \AttributeTok{main =} \StringTok{"Velocidade"}\NormalTok{)}

\CommentTok{\# Boxplot da Distância}
\FunctionTok{boxplot}\NormalTok{(cars}\SpecialCharTok{$}\NormalTok{dist,}
        \AttributeTok{main =} \StringTok{"Distância"}\NormalTok{)}
\end{Highlighting}
\end{Shaded}

\pandocbounded{\includegraphics[keepaspectratio]{classwork-1_files/figure-latex/Analise exploratória-1.pdf}}

\begin{Shaded}
\begin{Highlighting}[]
\FunctionTok{par}\NormalTok{(}\AttributeTok{mfrow =} \FunctionTok{c}\NormalTok{(}\DecValTok{1}\NormalTok{,}\DecValTok{1}\NormalTok{))}
\CommentTok{\#análise do gráfico}
\FunctionTok{plot}\NormalTok{(distancia, velocidade, }\AttributeTok{main =} \StringTok{"Hubble: Distância x Velocidade"}\NormalTok{)}
\end{Highlighting}
\end{Shaded}

\pandocbounded{\includegraphics[keepaspectratio]{classwork-1_files/figure-latex/Analise exploratória-2.pdf}}
O gráfico acima relaciona a velocidade das galáxias (eixo y) com a
distância entre elas (eixo x). Como pode-se observar, os pontos no
gráfico se concentram em uma linear crescente no meio do gráfico. Ou
seja, nota-se que, a medida que as distâncias aumentam, as velocidades
também aumentam.

\begin{Shaded}
\begin{Highlighting}[]
\NormalTok{coef }\OtherTok{\textless{}{-}} \FunctionTok{lm}\NormalTok{(velocidade}\SpecialCharTok{\textasciitilde{}}\NormalTok{distancia, }\AttributeTok{data =}\NormalTok{ dados)}
\NormalTok{coef}
\end{Highlighting}
\end{Shaded}

\begin{verbatim}
## 
## Call:
## lm(formula = velocidade ~ distancia, data = dados)
## 
## Coefficients:
## (Intercept)    distancia  
##      -48.34       505.84
\end{verbatim}

\begin{Shaded}
\begin{Highlighting}[]
\FunctionTok{plot}\NormalTok{(distancia, velocidade, }\AttributeTok{main =} \StringTok{"Regressão linear: Distância x Velocidade"}\NormalTok{)}
\FunctionTok{abline}\NormalTok{(coef, }\AttributeTok{col =} \StringTok{"blue"}\NormalTok{, }\AttributeTok{lwd =} \DecValTok{2}\NormalTok{)}
\end{Highlighting}
\end{Shaded}

\pandocbounded{\includegraphics[keepaspectratio]{classwork-1_files/figure-latex/Regressão Linear-1.pdf}}

Como dito anteriormente, existe relação linear entre a distância (\(r\))
e a velocidade (\$\upsilon \$), pois enquanto a distância aumenta, a
velocidade aumenta seguindo a seguinte relação:

\[\upsilon   = \beta_1 + \beta_2r + \varepsilon  \]
\[\upsilon  = -48.34  + 505.84r+ \varepsilon  \]

\section{Conclusão}\label{conclusuxe3o}

\end{document}
